\chapter*{Abstract}
\addcontentsline{toc}{chapter}{Abstract}

Rapid infrastructure expansion, urban gigaprojects, and agricultural initiatives across the Kingdom of Saudi Arabia under Vision 2030 require automated, scalable, and continuous environmental monitoring. However, satellite-based change detection in hyper-arid environments faces severe scientific challenges: high seasonal sun angle variations, soil moisture shifts, sabkha salinity, and spectral similarities between urban concrete and natural limestone deserts frequently trigger catastrophic false alert rates in conventional algorithms and untuned computer vision models.

This capstone project presents \textbf{TerraShift AI}, a production-ready, full-stack geospatial platform designed to identify, locate, measure, and analyze authentic land-surface changes using free Copernicus Sentinel-2 Level-2A multispectral imagery and Google AlphaEarth annual satellite foundation-model embeddings. To resolve the central research question—\textit{whether deep learning and foundation models genuinely outperform conventional remote sensing methods across diverse geographic environments}—we conduct a controlled, head-to-head empirical evaluation across four independently annotated Saudi reference corridors totaling $91.19\text{ km}^2$: (1) Riyadh ROSHN Sedra (residential earthworks on limestone desert), (2) Jeddah North Obhur (coastal infill and sabkha), (3) Al-Jouf Wadi As-Sirhan (center-pivot agricultural irrigation), and (4) an entirely held-out test corridor in Al-Kharj (logistics warehouse development and oasis activation).

Methodologically, we formulate a strict greedy $1$-to-$1$ bipartite event-matching protocol that strictly distinguishes true positive alerts from duplicate fragmented alerts and background false alarms, auditing negative control desert zones for zero-change integrity. We evaluate seven distinct model paradigms: (1) Untuned Siamese U-Net (RGB), (2) Tuned Siamese U-Net coupled with directional physical spectral guards ($\Delta\text{NDBI} \ge +0.04$, $|\Delta\text{NDVI}| \ge 0.10$), (3) Pure conventional multispectral index thresholding, (4) AlphaEarth 64-dimensional latent cosine dissimilarity ($\tau=0.12$), (5) AlphaEarth combined with Dynamic World land-use transition analysis, (6) AlphaEarth lightweight supervised vector difference classification, and (7) a hybrid AlphaEarth foundation model gated by physical electromagnetic rules.

Empirical results demonstrate that the Tuned Siamese model with physical spectral guards achieves the highest overall change detection accuracy with an Event F1 score of \textbf{61.5\%}, a Strict Precision of \textbf{45.5\%} (reaching 80.0\% in Jeddah and 55.6\% in Riyadh), and a Recall of \textbf{95.0\%}, while slashing false positive alerts by \textbf{69.6\%} (from 115 to 35) and completely eliminating false alarms in negative control zones ($3 \to 0$). In contrast, standalone AlphaEarth annual embeddings achieve \textbf{100.0\% strict precision} with zero false alarms across all corridors, but exhibit limited recall (21.2\%) on granular parcel-level developments without local adaptation. Pure physical multispectral indices achieve 95.0\% recall but suffer from elevated background noise (109.7 alerts/$100\text{ km}^2$, F1: 26.5\%), whereas untuned RGB models exhibit substantial false alarm rates (126.1 alerts/$100\text{ km}^2$).

Architecturally, TerraShift AI is implemented as a modern web platform combining Next.js 16 (React 19, MapLibre GL, TypeScript strict mode) with a high-performance Python FastAPI inference engine. The entire pipeline executes locally on consumer-grade CPU hardware with a mean pure compute latency of \textbf{0.495 seconds} per corridor and zero cloud GPU operating cost. The project establishes an open, scientifically validated benchmark for AI-assisted remote sensing in arid regions, providing an actionable tool for national geospatial compliance, municipal planning, and environmental conservation in the Kingdom.

\vspace*{1.0cm}
\noindent
\textbf{Keywords:} Remote Sensing, Change Detection, Satellite Foundation Models, Google AlphaEarth, Siamese Neural Networks, Sentinel-2, Hyper-Arid Monitoring, Saudi Vision 2030, Geospatial AI.
\newpage
